\section{Introduction}

Missing covariate data is a pervasive problem in observational studies, and the standard linear modelling tools used for differential expression analysis cannot be fit to records that are incomplete. The simplest remedy, complete case (CC) analysis, discards any individual with at least one missing value; this practice reduces statistical power and, depending on the missingness mechanism, can introduce substantial bias into model coefficients \cite{sterne2009multiple,little2019statistical}. A second widely used strategy, single imputation (SI), replaces each missing value with the most likely prediction from a supplementary model. While SI preserves the full sample, it treats imputed values as if they were observed and therefore underestimates standard errors and inflates confidence in the downstream estimates \cite{rubin1987multiple,sterne2009multiple}. Multiple imputation (MI) was developed specifically to propagate the uncertainty associated with filling in missing values: $M$ imputed datasets are created, each is analysed separately, and the resulting estimates and standard errors are combined using Rubin's rules \cite{rubin1987multiple,marshall2009combining}. Barnard and Rubin's small-sample adjustment to the pooled degrees of freedom further stabilises inference when the completed fraction is modest \cite{barnard1999smalls}.

Despite decades of methodological development in epidemiology, there has been no concerted effort to determine the most advantageous way of handling missing covariate data in transcriptomic studies. Historically, a large fraction of RNA-sequencing (RNA-seq) work has been performed in in-vitro or in-vivo models that do not suffer from missing covariates, and experimental designs with a small number of controlled variables often avoid the problem altogether. The calculus is changing quickly. As sequencing costs continue to fall, transcriptomic assays are being deployed at scale inside human observational cohorts such as ECHO-PATHWAYS \cite{paulson2022lewinn,paquette2023placental}, where missing covariate data is the rule rather than the exception. Without tailored methods, analysts revert to CC or SI, each with well-understood statistical hazards.

A direct application of MI to RNA-seq is not straightforward. The most effective MI predictor matrices include the outcome variable, which helps preserve the outcome--covariate relationship in the imputed values \cite{vanbuuren2011mice}. For RNA-seq experiments with tens of thousands of genes, however, placing every gene in the predictor matrix would require an equally large number of participants and is typically infeasible. An obvious alternative --- constructing one set of $M$ imputed datasets per gene --- quickly becomes computationally prohibitive. Analogous challenges arose in epigenome-wide association studies (EWAS), where methods based on grouping CpG sites together in the predictor matrix were introduced to recover tractability \cite{mill2013epigenomewide,weinhold2020statistical}. Inspired by these ideas, we developed RNAseqCovarImpute, an R package that bins genes into small groups, constructs $M$ imputed datasets separately within each bin, and pools results with Rubin's rules. The package is fully compatible with the limma-voom pipeline \cite{law2014voom,ritchie2015limma}, which is the de facto standard for RNA-seq differential expression. We assess RNAseqCovarImpute in three simulation studies and a real-world application to placental transcriptomics.

\section{Related work}

\paragraph{RNA-seq differential expression.} The limma-voom framework estimates gene-wise precision weights from the mean-variance trend of log counts per million and uses those weights inside limma's moderated linear model \cite{law2014voom,ritchie2015limma}. Alternative pipelines such as edgeR \cite{robinson2010edger} and DESeq2 \cite{love2014deseq2} fit negative binomial models directly to integer counts. Normalisation by trimmed mean of M-values (TMM) \cite{robinson2010tmm} is standard across all three pipelines. Downstream enrichment analyses such as camera \cite{wu2012camera} and GAGE \cite{luo2009gage} operate on gene-level statistics or rankings produced by these methods, which matters for compatibility with MI approaches discussed below.

\paragraph{Imputation in the omics context.} Single-cell RNA-seq has motivated a large body of work on imputing missing or zero counts to address dropout \cite{hao2024scrnaseq}. These methods target sparsity in the outcome matrix rather than missingness in covariates, and they do not address the problem tackled here. Generic imputation tools such as mice (chained equations for mixed-type data) \cite{vanbuuren2011mice} and missForest (nonparametric imputation via random forests) \cite{stekhoven2012missforest} work at the covariate level but were not designed for outcome dimensionalities exceeding sample size.

\paragraph{Missing data in EWAS.} Previous work on multiple imputation for epigenome-wide association studies binned CpG sites to make outcome-informed MI tractable \cite{mill2013epigenomewide,weinhold2020statistical}. Reported tradeoffs favoured MI for sensitivity but at the cost of higher FDR in some settings. Methylation arrays such as the Illumina MethylationEPIC BeadChip measure around $850{,}000$ sites \cite{pidsley2016critical}, a dimensionality that is an order of magnitude larger than a typical RNA-seq experiment, and so translating EWAS MI methods directly to transcriptomics is not necessarily a well-posed problem.

\paragraph{Simulation of RNA-seq signal.} We evaluate our method using both real ECHO-PATHWAYS counts \cite{paulson2022lewinn} and synthetic counts generated with seqgendiff \cite{gerard2020data}, which uses binomial thinning of a real count matrix to induce user-specified differential expression coefficients while preserving realistic variance structure.

\paragraph{FDR control and missing data methodology.} False discovery rate control via the Benjamini--Hochberg procedure \cite{benjamini1995controlling} is the standard multiple-comparison correction for high-dimensional assays. Broader methodology for missing data in epidemiology is reviewed in \cite{sterne2009multiple,little2019statistical}, and pooling procedures after MI are summarised in \cite{marshall2009combining}. To our knowledge, no prior work has combined these tools into a complete differential expression pipeline that accommodates missing covariates in RNA-seq.
